\section{Background: what the four questions rest on}
\label{sec:background}

This report is written to be read by someone with a first course in statistics, a working idea of
what fine-tuning a small language model means, and a lay familiarity with mortgage lending. Everything
technical the four questions rely on is defined once here, from the ground up, with the exact equations and
a worked example wherever a number would otherwise be mysterious. Nothing in this section reports a
result; it is the vocabulary. Almost every subtlety below exists to make one comparison honest --- a
tuned guard against the \emph{same model before tuning}, on identical rows, split into sources the
guard trained on and sources it never saw --- and \Cref{fig:design} previews that design on one page.

\begin{figure}[t]\centering
\includegraphics[width=0.70\linewidth]{experiment_design.png}
\caption{\textbf{The study at a glance (Acts~I--II).} From one fixed panel and one frozen training
manifest, every guard we compare --- the untuned \textbf{base}, the \textbf{SFT-tuned} guard, and the
retraining-free \textbf{composition} --- is read by the \emph{same} single-token scorer on the
\emph{same} represented and transfer test rows, and each is compared only against its \emph{own} base.
That paired difference $\Delta$, aggregated by a family-aware bootstrap, is the estimand. Act~III then adds
the mortgage and finance/health/law domain evaluations on top of this same machinery.}
\label{fig:design}
\end{figure}

\subsection{The guard: a single-token logit-difference head}
\label{sec:bg-guard}

We do not ask the guard to write a paragraph explaining its decision. We ask it for exactly one
word and look only at how strongly it leans toward that word. Concretely, a prompt $x$ is wrapped in a
fixed instruction template asking for a one-word verdict, and we read the model's raw output scores
--- its \emph{logits} --- at the final position for the two verdict words \code{safe} and
\code{unsafe}. A logit is the model's un-normalized, pre-probability score for a candidate next token:
larger means the model favors that token. Writing $z_{\text{unsafe}}(x,t)$ and $z_{\text{safe}}(x,t)$
for those two logits at the final prompt position $t$, the guard's stored score is their difference,
\begin{equation}
\label{eq:score}
s(x) \;=\; z_{\text{unsafe}}(x,t) \;-\; z_{\text{safe}}(x,t),
\end{equation}
which is positive when the model leans \code{unsafe}, negative when it leans \code{safe}, and near zero
when it is undecided. To read $s(x)$ as a probability we pass the same two logits through a two-way
softmax --- giving, strictly, the probability of \code{unsafe} \emph{conditional on the verdict being one
of the two tokens} $\{$\code{safe}$,$\code{unsafe}$\}$, not the model's full next-token distribution:
\begin{equation}
\label{eq:prob}
p(\text{unsafe}\mid x) \;=\; \frac{\exp\!\big(z_{\text{unsafe}}\big)}{\exp\!\big(z_{\text{safe}}\big)+\exp\!\big(z_{\text{unsafe}}\big)}.
\end{equation}
For example, if $z_{\text{unsafe}}=2.0$ and $z_{\text{safe}}=1.0$, then $s(x)=1.0$ and
$p(\text{unsafe}\mid x)=e^{2}/(e^{2}+e^{1})\approx 0.73$. This is what we mean by a
``single-token logit-difference head'': one forward pass, two numbers, one score. The raw logits are
the canonical stored value; probabilities in \Cref{eq:prob} and any temperature-rescaled version are
\emph{derived} from them. Reading the model's own generated verdict text is deliberately never used as
the primary score, because a free-text answer is harder to score consistently and hides the model's
actual confidence.

\subsection{Base, SFT, and LoRA}
\label{sec:bg-sft}

The \emph{base} is the general instruction-tuned chat model \emph{before} we turn it into a guard. It
is not untrained --- it can already follow the ``answer \code{safe} or \code{unsafe}'' instruction
zero-shot --- so ``base'' throughout means ``before the guard fine-tune,'' not ``blank.''
\emph{Supervised fine-tuning} (SFT) then trains that model on labeled examples: for each training
prompt we know the correct verdict, and we nudge the model to raise the probability it assigns to that
verdict word. Rather than update all of the model's weights, we use \emph{LoRA} (Low-Rank Adaptation)
\citep{hu2021lora}: we freeze the original weights and insert a small number of extra trainable
parameters (a low-rank ``adapter'') into selected weight matrices. LoRA is cheap, fast, and reversible
--- you can ship the tiny adapter on top of the frozen base --- and it is the standard way guards are
built in practice \citep{elesedy2024loraguard}. The precise adapter size and which matrices it touches
are part of the frozen recipe in \Cref{sec:setup}.

\subsection{Average precision, macro-AP, and the accuracy trap}
\label{sec:bg-ap}

Because $s(x)$ is a continuous score, we can grade a guard \emph{without} committing to a cutoff, by
asking how well it \emph{ranks}. \emph{Average precision} (AP) collapses the whole
precision--recall curve into a single number in $[0,1]$: it is near $1$ when the guard tends to give
genuinely unsafe prompts higher scores than safe ones, and it needs no threshold. Formally, sweeping
the cutoff from high to low, AP is the recall-weighted average of the precision achieved at each
true-positive, $\text{AP}=\sum_n (R_n-R_{n-1})\,P_n$, with $P_n,R_n$ the precision and recall after the
$n$-th ranked item; we use the tie-aware, non-interpolated \code{scikit-learn} definition so the number
is exactly reproducible.

\begin{background}{A worked AP example}
Rank five prompts by their guard score, highest first, and write their true labels:
$[\text{unsafe},\ \text{unsafe},\ \text{safe},\ \text{unsafe},\ \text{safe}]$. There are three unsafe
prompts, so each one contributes recall $\tfrac13$. The precision at each unsafe prompt as we go down
the list is $\tfrac11{=}1.00$ (rank 1), $\tfrac22{=}1.00$ (rank 2), and $\tfrac34{=}0.75$ (rank 4). So
$\text{AP}=\tfrac13(1.00)+\tfrac13(1.00)+\tfrac13(0.75)\approx 0.917$. A perfect ranking (all three
unsafe prompts on top) would score $1.00$; ranking the safe prompts first would score much lower.
\end{background}

\emph{Macro}-AP computes AP separately on each benchmark and then averages those per-benchmark values
with \emph{equal weight}, so one large benchmark cannot dominate a handful of small ones; this is the
primary metric behind every one of the four questions.

\begin{takeaway}{Why not just report accuracy?}
On a test set that is $95\%$ safe, a guard that blindly answers \code{safe} to everything scores $95\%$
accuracy while catching \emph{zero} attacks. Accuracy rewards guessing the majority class; AP and
macro-AP instead reward putting the unsafe prompts on top, which is the behavior a guard is for. This
is why ranking, not accuracy, is the currency of this report.
\end{takeaway}

\subsection{Two evaluation regimes: represented-source vs.\ dataset-held-out transfer}
\label{sec:bg-regimes}

The single most important distinction in this report is \emph{where the test data came from}.
\textbf{Represented-source} test rows are held-back examples drawn from a dataset the guard
\emph{trained on} (different rows, same source). \textbf{Dataset-held-out transfer} rows come from
datasets whose examples were \emph{never used in training at all}. A gain that appears on
represented-source tests but not on transfer tests is precisely the signature of a guard
\emph{specializing} to its training sources rather than becoming a broadly better moderator. One
honesty caveat travels with the word ``held out'': it means \emph{dataset-held-out} (those rows were
withheld from fine-tuning), \emph{not} sealed away from the researchers during development --- a
distinction that forces the retrospective framing discussed in \Cref{sec:bg-estimand}.

\subsection{From scores to decisions: calibration, operating point, TPR/FPR, macro vs.\ pooled}
\label{sec:bg-op}

Ranking is threshold-free, but a deployed guard must eventually say \code{safe} or \code{unsafe}. An
\emph{operating point} is the single cutoff on $s(x)$ that converts scores into hard decisions. Picking
it trades two error rates against each other: the \emph{true-positive rate} (TPR, or recall) is the
fraction of genuinely unsafe prompts the guard flags, and the \emph{false-positive rate} (FPR) is the
fraction of genuinely safe prompts it wrongly flags. Raising the cutoff lowers both; lowering it raises
both. Before thresholding we \emph{calibrate}: we fit a single positive \emph{temperature} that rescales
the logits (temperature scaling \citep{guo2017calibration}) so the derived probabilities in
\Cref{eq:prob} are better behaved. Crucially, both the temperature and the cutoff are fit only on a
separate held-back \emph{calibration} split --- never on the transfer or stress data we report on --- so
the threshold is not quietly tuned to the test.

Two error rates can be averaged two ways, and the choice matters. \emph{Macro} rates compute the rate
on each benchmark first and then average across benchmarks (equal weight per benchmark). \emph{Pooled}
rates lump every row together first and then compute one rate (equal weight per row, so large
benchmarks dominate). We report both because they can disagree.

Keep the two questions apart: AP answers ``does it rank?'' and the operating point answers ``is there a
usable decision boundary?'' A guard can rank well and still have no cutoff that survives off-source,
because the score \emph{distribution} shifts even when the \emph{ordering} holds --- a gap Acts~I and~II
both measure (\Cref{sec:actI-op,sec:actIII-operating}).

\subsection{Estimands, the fixed panel, and the paired hierarchical bootstrap}
\label{sec:bg-estimand}

\begin{background}{Three terms that appear in the verdicts, in plain words}
\textbf{Lower confidence bound (LCB).} A one-sided version of an interval. Saying ``the gain is
\AdaHGain{}, LCB \AdaHGainLCB'' --- the actual RQ1 result of \Cref{sec:adaptation} --- means: the
estimate is \AdaHGain{}, and under resampling the value stayed above \AdaHGainLCB{} in $97.5\%$ of the
redraws. A criterion of the form ``LCB $>0$'' is therefore a demand that even
the pessimistic end of the range still be a gain. (UCB is the mirror image, used for quantities we want
to bound from \emph{above}, such as a loss.) This example is the generated macro rather than a typed
number, because an earlier revision left it quoting a superseded panel's value.

\textbf{Non-inferiority margin.} How much you are willing to lose on one axis to win on another, fixed
\emph{before} looking. A margin of $-0.02$ says: a drop of up to two AP points is acceptable, more is
not. This is stricter than ``did it get worse?'' --- it must be provably \emph{not much} worse, so a
noisy result fails the test rather than passing it by default. \Cref{sec:adaptation} contains the one
place in this report where a preregistered criterion of this kind is failed and reported as failed.

\textbf{Bonferroni split.} When you test two things at once, each gets half the error budget
($\alpha=0.05$ becomes $0.025$ each), so asking two questions cannot double your chance of a lucky
answer. It is the cheapest, most conservative way to pay for multiplicity.
\end{background}

An \emph{estimand} is the exact quantity our numbers describe. Here it is the average
before-versus-after change \emph{for the four specific checkpoints we study}, not for ``small guards''
in general. Those four checkpoints are a \emph{fixed, purposively chosen panel} --- picked deliberately,
not sampled at random from a population --- so we do not attach uncertainty to the choice of models and
we do not generalize to unnamed architectures. Every comparison is \emph{paired}: an adapter is always
compared against \emph{its own base} on the \emph{same rows}, so nothing but the fine-tune differs.

To put a $95\%$ range around a paired change we use a \emph{paired hierarchical bootstrap}. A bootstrap
estimates how much a result would wobble under resampling by rebuilding the data many times from itself
and recomputing; the spread of the recomputed values becomes the interval. Ours is \emph{hierarchical}
because it resamples at two levels at once --- which fine-tuning \emph{seeds} were drawn (within each
fixed checkpoint), and which groups of near-duplicate evaluation rows (``families,''
\Cref{sec:setup-decon}) are counted, via one Poisson(1) re-count weight per family. It is \emph{paired}
because each adapter stays yoked to its base. Because the training rows and their order are held fixed,
seed-to-seed variation reflects initialization, dropout, and execution nondeterminism --- not
which examples the guard happened to see. The resulting intervals are \emph{conditional} on this panel,
these datasets, and this family graph; they are descriptive, not significance tests, and not claims
about a wider population. This is why Acts~I--II report point estimates, intervals, and leave-one-out
sensitivity checks but no accept/reject test: they were analyzed \emph{after} the benchmarks had been
inspected during development, and clean provenance does not make an inspected benchmark prospective
(see the scope note on p.~1 and \Cref{sec:honesty}).

\subsection{Fine-tuning the guard: LoRA-SFT}
\label{sec:bg-objectives}

We fine-tune each guard with \textbf{SFT} (supervised fine-tuning): a cross-entropy loss that directly
raises the model's log-probability of the single correct verdict token, applied as a parameter-efficient
LoRA adapter (the exact recipe is \Cref{tab:panel-recipe}). SFT is the one training recipe used throughout;
Act~I measures what it changes, and Act~II asks whether we can recover what it gives up \emph{without}
any further tuning.

\subsection{Output-space composition}
\label{sec:bg-composition}

Act~II introduces a retraining-free remedy. \emph{Output-space composition} combines two guards by
averaging their calibrated \emph{output scores} --- not their internal weights. It is portable, because
it needs only two comparable scores rather than interpolable parameters, but it costs two inference
passes because both models must run. This contrasts with \emph{weight-space} methods such as WiSE-FT and
model soups \citep{wortsman2022wiseft,wortsman2022modelsoups}, which average the models' weights and so
require the models to be weight-compatible. The exact composition rule is given later as
\Cref{eq:comp}.

\subsection{The mortgage/HMDA vocabulary and the dual-label idea}
\label{sec:bg-mortgage}

Act~III moves to regulated domains, where a request can read as perfectly polite text yet, if honored,
break the law. The mortgage vocabulary below is what the benchmark encodes.

\begin{background}{Mortgage terms, in plain words}
\textbf{HMDA} --- the U.S.\ Home Mortgage Disclosure Act \citep{hmda2022}: lenders publicly report
loan-level records (purpose, amount band, action taken, denial reason), which we use only as
\emph{de-identified, banded fact sheets}, never verbatim. \textbf{Fair lending} bars credit
discrimination on protected traits. \textbf{Redlining} is refusing or worsening credit across a whole
neighborhood as a proxy for race. \textbf{Disparate treatment} is intentional differential treatment;
\textbf{disparate impact} is a facially neutral rule that falls harder on a protected group.
\textbf{Proxy (coded) discrimination} uses a stand-in --- ZIP code, preferred language --- for a
protected class. \textbf{ATR/QM} is the ability-to-repay / qualified-mortgage rule.
\textbf{Adverse-action notice} is the legally required true reason for a denial. \textbf{UDAAP} covers
unfair, deceptive, or abusive acts and practices.
\end{background}

\noindent On top of that vocabulary the benchmark adds one construct
(\Cref{sec:actIV-mortgage-design}): two \emph{separately assigned} labels per request, general safety $G$ and
mortgage policy $D$, isolating the stratum a general-safety score cannot see.
